% !TEX root = ../main-es.tex
% common/meta-es.tex — metadatos en español emitidos como comandos de mitthesis.cls
% Las strings condicionales al español se registran aquí (único sitio).

% ===== Strings de la clase al español =====
\SuppressMonthError                      % Agosto no es mes MIT válido
\SupervisorDesignation{Director}         % default: "Thesis Supervisor"
\ThesisReaderName{Jurado}                % default: "Thesis Reader"
\ThesisCommitteeName{Comité de Tesis}    % default: "THESIS COMMITTEE"
\renewcommand{\abstractname}{Resumen}
\renewcommand{\contentsname}{Índice general}
\renewcommand{\listfigurename}{Índice de figuras}
\renewcommand{\listtablename}{Índice de tablas}
\renewcommand{\refname}{Referencias}
\renewcommand{\bibname}{Referencias}
\renewcommand{\indexname}{Índice alfabético}
\renewcommand{\figurename}{Figura}
\renewcommand{\tablename}{Tabla}
\renewcommand{\chaptername}{Capítulo}
\renewcommand{\appendixname}{Apéndice}

% ===== Datos de la tesis (comandos MIT) =====
\title{Semántica categórica en un DSL de instrumentación agéntica\\%
      \large Composicionalidad como garantía estructural frente a la degradación en profundidad\\\normalsize
      % El caso de \OASys
      }

\Author{Over Haider Castrillón Valencia}{Maestría en Matemáticas Aplicadas y Ciencias de la Computación}

\Degree{Maestría en Matemáticas Aplicadas y Ciencias de la Computación}%
       {Escuela de Ingeniería, Ciencia y Tecnología}

\Institution{Universidad del Rosario}

% Logo institucional en la portada (parche UR en mitthesis.cls)
\ThesisLogo{ur-logo}[3.2cm]

\Supervisor{Vivian Estefanía Beltrán Parada}{Profesor del Programa}%
            [Maestría en Matemáticas Aplicadas y Ciencias de la Computación]

% Acceptor mapeado a Coordinador del Programa (figura UR equivalente)
\Acceptor{Daniel Alfonso Bojaca Torres}{Coordinador del Programa}%
         {Maestría en Matemáticas Aplicadas y Ciencias de la Computación}

% Tres jurados (placeholders mientras se asigna el comité)
\Reader{(Jurado 1)}{Profesor del Programa}{Maestría en MACC}
\Reader{(Jurado 2)}{Profesor del Programa}{Maestría en MACC}
\Reader{(Jurado 3)}{Profesor del Programa}{Maestría en MACC}

\DegreeDate{August}{2026}   % mes MIT válido (SuppressMonthError silencia otros)
\ThesisDate{Agosto de 2026}

% ===== Atajos del proyecto para capítulos y apéndices =====
% ( Conservamos \tesisresumen y \tesispalabras para uso del abstract;
%   el resto ya vive en los comandos MIT de arriba. )
\newcommand{\tesispalabras}{Orquestación de agentes; teoría de categorías; categorías monoidales; semántica denotacional; verificación estática; \OASys}
\newcommand{\tesisresumen}{%
  Este trabajo da semántica denotacional al fragmento núcleo del lenguaje \dsl{.os} de
  \OASys{} como morfismos de una categoría de Freyd sobre una categoría monoidal simétrica.
  La separación entre un plano de control determinista y un plano de efecto estocástico
  (el LLM), ya realizada por el motor en ADR-003, se formaliza como una categoría de Freyd.
  Se enuncia y demuestra un teorema de no-interferencia composicional para el sub-fragmento
  fork-join, y se implementa un verificador estático de referencias que convierte una clase
  de fallos de ejecución en fallos de compilación. Se mide el efecto sobre la degradación
  por profundidad frente a un baseline externo y a una ablación del verificador sobre el
  mismo motor.%
}

% ===== Strings del frontmatter =====
\newcommand{\tesislabelkeyw}{Palabras clave}
\newcommand{\tesislabeldecl}{Declaración de autonomía intelectual.}
\newcommand{\tesisdeclcorta}{Versión corta --- el texto firmado íntegro está en el Apéndice \ref{app:declaraciones}.}
